\chapter{Introduction}
\label{ch:introduction}

The study of transverse momentum ($p_T$) spectra of particles produced in high-energy proton-proton and heavy-ion collisions provides critical insights into the thermodynamic properties and collective dynamics of the strongly interacting medium. A fundamental challenge in modeling these spectra lies in the tension between two distinct physical regimes: the soft "bulk" particle production at low $p_T$ (dominated by thermal emission and collective hydrodynamic expansion) and the hard scattering tail at high $p_T$ (dominated by perturbative quantum chromodynamics and jet fragmentation). 

Historically, phenomenological models such as the Boltzmann-Gibbs Blast-Wave (BGBW) framework and the thermodynamically consistent Tsallis distribution have been employed to reconcile these regimes, extracting kinetic freeze-out temperatures and transverse flow velocities. However, attempts to unify these regimes have led to the proposition of increasingly complex multi-component models. 

This thesis investigates one such proposal: a three-component model aiming to describe the full $p_T$ spectrum using multiple moving J\"uttner (relativistic Maxwell-Boltzmann) distributions. Evaluating the validity of such highly parameterized models requires moving beyond simple visual fits to rigorous statistical and physical verification.

To this end, we introduce a novel methodological contribution: an autonomous, multi-agent artificial intelligence framework designed to systematically validate phenomenological physics models against both experimental data and literature consensus. By applying this AI-augmented framework to the proposed J\"uttner model, we execute an end-to-end validation pipeline to definitively assess its statistical rigor, physical justification, and structural stability against established literature baselines.
